\documentclass[11pt]{article}
\usepackage{fontspec}
\usepackage[french]{babel}
\usepackage{amsmath}
\usepackage{amssymb}
\usepackage{amsthm}
\usepackage{mathrsfs}
\usepackage{tikz-cd}
\usetikzlibrary{babel}
\usepackage{enumitem}

\newcommand{\printedpage}[1]{\marginpar{\footnotesize p.~#1}}
\newcommand{\msp}[1]{\marginpar{\footnotesize [#1]}}
% printed notation
\newcommand{\cl}{^{\tilde{}}}% tilde printed as exponent after a symbol or a parenthesis
\newcommand{\defeq}{\overset{\text{déf}}{=}}
\newcommand{\encircled}[1]{\tikz[baseline=(X.base)]\node[draw,circle,inner sep=1pt](X){#1};}
\newcommand{\encoval}[1]{\tikz[baseline=(X.base)]\node[draw,rounded corners=4pt,inner sep=2pt](X){#1};}

\begin{document}

\printedpage{44}
% head: A. Grothendieck : Dérivateurs, Ch. XVII, Édition des Universités de Montpellier II et Paris VII
\[
\begin{tikzcd}
TC \arrow[r, hook] \arrow[d, leftrightarrow] & C \arrow[d, leftrightarrow] \\
F & TF \arrow[l, hook'] \ ,
\end{tikzcd}
\tag*{(5.2)}
\]
et $TC = C \cap W$, le fait que $(W, C, F)$ est un triple de Quillen est le corollaire 4 p. 36 de XV : Les conditions $TC = C \cap W$, $TF = F \cap W$, $W = TF \circ TW$ sont équivalentes, et équivalentes au fait que $(W, C, F)$ soit un triple de Quillen.

Par la définition de Quillen des triples \textit{clos,} par les propriétés
\[
\underbrace{W = TF \circ TC}_{\substack{\text{vraie pour tous}\\ \text{le triples}\\ \text{de Quillen}}}, \qquad C = (TF)^*, \qquad \underbrace{F = (TC)_*}_{\substack{\text{vraie dans}\\ \text{le cas actuel}}}\ , % [sic] le triples -> les triples, agreement with "tous"
\]
il reste à exprimer que $C = (TF)^*$, ce qui équivaut à $C$ stable par facteurs directs (cf. loc. cit. cor. 1, p. 31), q.e.d.

\textbf{Remarque 5.4.} Les énoncés 5.2 et 5.3, et la relation $T\tilde{C} \subset (TC)^\natural$ dans proposition 5.1, ne font pas appel à la totalité de l'hypothèse que $(\mathcal{M}, W, C)$ soit une catégorie à cofibrations spéciale, mais seulement que l'on ait une catégorie à cofibrations, avec $C \mathrel{\leftarrow\!\!\!+\!\!\!\rightarrow} \underbrace{TF}_{\defeq C_*} \subset W$, % ?? uncertain: arrow symbol between C and TF -- double-headed arrow crossed by a small bar/plus sign
\msp{page 89}
i.e. une catégorie avec $W$, $C$ satisfaisant Cof$_0$, Cof$_1$, Cof$_2$ de la page 39, plus $TF \defeq C_* \subset W$, plus l'axiome de factorisation pour $(C, TF)$ (lequel implique Cof$_3$).

\textbf{Corollaire 5.5.} \textit{Soit $(\mathcal{M}, W, C)$ une catégorie à cofibrations spéciale, telle qu'il existe une} petite \textit{partie $TC_0 \subset TC = C \cap W$, avec}
\[
TC \ \subset \ TC_0\cl\ ,
\tag*{(5.5.1)}
\]
\textit{ou ce qui revient au même, $(TC)\cl = (TC_0)\cl$, i.e.}
\[
\underbrace{F}_{\defeq (TC)_*} \ = \ (TC_0)_*\ .
\tag*{(5.5.2)}
\]
\textit{Alors $(W, C, F)$ est un triple de Quillen. Ce triple est clos si et seulement si $C$ est stable par facteurs directs.}

En effet, le théorème de factorisation dans XV §1 implique que moyennant l'hypothèse faite sur $TC_0 \subset TC$, le couple $(TC, F)$ satisfait à la condition de factorisation. En effet, par les conditions b, b$'$ de la définition 4.4, $TC$ a les propriétés de stabilité voulues (contient les isomorphismes, est stable par cochangement de base et par composition transfinie).

\msp{page 90}
\textbf{Définition 5.6.} Les triples de Quillen déduits comme dans 5.5 d'une catégorie à cofibrations spéciale sont dits \textit{spéciaux.} Ainsi un triple de Quillen $(W, C, F)$ dans $\mathcal{M}$ est spécial s'il satisfait les conditions a) à d) dans la définition (4.4), concernant $(W, C)$ exclusivement, plus les deux conditions suivantes.
% foot: Version 9 février 2019

\printedpage{45}
% head: A. Grothendieck : Dérivateurs, Ch. XVII, Édition des Universités de Montpellier II et Paris VII
\begin{itemize}[leftmargin=3em]
\item[\encircled{c$'$}] $\exists$ petite partie $TC_0$ de $TC$, telle que $TC \subset (TC_0)\cl$.
\item[\encircled{e}] $F = (TC)_*$.
\end{itemize}

\textbf{Corollaire 5.7.} \textit{Soit $(W, C, F)$ un triple de Quillen} clos \textit{dans $\mathcal{M}$. Pour que ce triple soit spécial, il faut et il suffit qu'il satisfasse les conditions suivantes.}
\begin{itemize}[leftmargin=3em]
\item[\encircled{a}] \textit{$\mathcal{M}$ stable par petites $\varinjlim$\,, et ses éléments sont accessibles.}
\item[\encoval{c, c$'$}] \textit{Il existe une petite partie $C_0$ de $C$, $TC_0$ de $TC$, telles que l'on ait}
\[
C \ = \ C_0\cl\ , \qquad TC \ = \ (TC_0)\cl\ .
\tag*{(5.7.1)}
\]
\end{itemize}
En effet, toutes les autres conditions a) à e), y inclus c$'$) de 5.6, sont automatiquement satisfaites.

\msp{page 91}
\textbf{Proposition 5.8.} \textit{Considérons un couple de foncteurs adjoints}
\[
\begin{tikzcd}
\mathcal{M} \arrow[r, shift left, "v_!"] & \mathcal{M}' \arrow[l, shift left, "v^*"] \ ,
\end{tikzcd}
\tag*{(5.8.1)}
\]
\textit{où $\mathcal{M}$ est munie d'un triple de Quillen spécial clos $(W, C, F = (TC)_*)$, avec $TC = C \cap W$} (cf. définition 5.6). \textit{On suppose satisfaites les conditions suivantes} (comparer prop. 4.6, relations (4.6.1), et (4.7.4)) :
\[
\left\{
\begin{array}{ll}
\text{a)} & \mathcal{M}' \ \textit{stable par petites}\ \varinjlim\ \textit{, ses éléments sont accessibles.}\\[1.5ex]
\text{b)} & v^* \ \textit{commute aux petites}\ \varinjlim\ .\\[1.5ex]
\text{c)} & v^*(v_!(C)) \subset C,\ v^*(v_!(TC)) \subset TC.
\end{array}
\right.
\tag*{(5.8.2)}
\]
\textit{Définissons $(W', C', F')$ dans $\mathcal{M}'$} (comparer prop. 4.6, notamment 4.6.2) \textit{par}
\[
\left\{
\begin{array}{lcl}
W' & = & (v^*)^{-1}(W)\\[1ex]
C' & = & v_!(C)\cl\\[1ex]
F' & = & (TC')_*\ , \qquad \textit{où}\ \ TC' \ \defeq \ C' \cap W'\ .
\end{array}
\right.
\tag*{(5.8.3)}
\]
\textit{Alors $(W', C', F')$ est un triple de Quillen clos spécial dans $\mathcal{M}'$, et on a de plus}
\[
F' \ = \ (v^*)^{-1}(F)\ , \qquad TF' \ = \ (v^*)^{-1}(TF)\ ,
\tag*{(5.8.4)}
\]
\[
TC' \ = \ (v_!(TC))\cl\ .
\tag*{(5.8.5)}
\]

\msp{page 92}
\textsc{Démonstration}. Par 4.6, on sait déjà que $(\mathcal{M}', W', C')$ est une catégorie à cofibrations spéciale. (\textbf{N.B.} Ce qui est noté $C'$ ici, était noté $\tilde{C}'$ dans la proposition 4.5. On doit utiliser la dernière assertion de 4.6, où on suppose $C$ stable par facteurs directs, OK car $(W, C, F)$ supposé clos.) Pour prouver que $(W', C', F' \defeq (TC')_*)$ est un triple de Quillen
% foot: Version 9 février 2019
\printedpage{46}
% head: A. Grothendieck : Dérivateurs, Ch. XVII, Édition des Universités de Montpellier II et Paris VII
spécial, il reste, par 5.3 et définition 5.6, à prouver qu'il existe une petite partie $TC'_0$ de $TC'$, telle que
\[
(\ TC'_0 \ \subset\ )\ TC' \ \subset \ (TC'_0)\cl\ .
\tag*{(5.8.6)}
\]
(Ce triple de Quillen sera de plus clos, puisque $C'$ est Q-clos à gauche par définition.) Soit donc $TC_0$ une petite partie de $TC$ telle que
\[
TC \ \subset \ (TC_0)\cl\ ,
\tag*{(5.8.7)}
\]
qui existe par hypothèse sur $(W, C, F)$. On prendra
\[
TC'_0 \ \defeq \ v_!(TC_0)\ .
\tag*{(5.8.8)}
\]
On a vu dans 4.7 qu'on a
\[
v_!(TC) \ \subset \ TC'
\tag*{(5.8.9)}
\]
(résulte de l'hypothèse $v^*(v_!(TC)) \subset TC$ dans (5.8.2. c)), a fortiori
\[
TC'_0 \ = \ v_!(TC_0) \ \subset \ TC'\ ,
\]
\msp{page 93}
il reste à prouver
\[
TC' \ \overset{?}{\subset} \ (TC'_0)\cl\ ,
\tag*{(5.8.10)}
\]
ce qui équivaut à $(TC')\cl \subset (TC'_0)\cl$, ou encore à
\[
(TC'_0)_* \ \overset{?}{\subset} \ (TC')_* \ \defeq \ F'\ ,
\tag*{(5.8.10 bis)}
\]
alors qu'on a a priori l'inclusion inverse (déduite de $TC'_0 \subset TC'$)
\[
F' \ \subset \ (TC'_0)_*\ .
\tag*{(5.8.11)}
\]
Or par adjonction, on a
\[
\begin{aligned}
(\ \underbrace{TC'_0}_{v_!(TC_0)}\ )_* \ &= \ (v_!(TC_0))_*\\
&= \ (v^*)^{-1}(\ \underbrace{(TC_0)_*}_{= (TC)_* \defeq F}\ )\\
&= \ (v^*)^{-1}(F)\ ,
\end{aligned}
\]
donc (5.8.11) s'écrit
\[
F' \ \subset \ (v^*)^{-1}(F)\ , \qquad \text{i.e.}\ \ v^*(F') \ \subset \ F
\tag*{(5.8.12)}
\]
(comparer (4.7.6)), et il faut montrer que cette inclusion est une égalité (sur quoi j'avais buté p. 80, 81 \ldots), ce qui n'est autre que la première des deux relations (5.8.4). Mais c'est sous la forme initiale (5.8.10) (transposée de (5.8.4)) que nous allons prouver la chose, finalement. Notons que l'on a
% foot: Version 9 février 2019
\printedpage{47}
% head: A. Grothendieck : Dérivateurs, Ch. XVII, Édition des Universités de Montpellier II et Paris VII
\msp{page 94}
\[
(\underbrace{v_!(TC_0)})\cl_{TC'_0} \ = \ (v_!(TC))\cl\ ,
\]
puisque $TC$ est la clôture de $TC_0$ pour les opérations de cochangement de base, de composition transfinie, et de facteurs directs, et que $v_!$ commute à ces opérations. Donc on doit montrer simplement
\[
TC' \ \subset \ (v_!(TC))\cl
\tag*{(5.8.12 bis)}
\]
(où il n'est plus question de $TC_0$). Mais le théorème de factorisation s'applique au couple
\[
\Big( (v_!(TC))\cl,\ (v_!(TC))_* = (v^*)^{-1}(F) \Big)\ ,
\]
puisque $(v_!(TC))\cl = (TC'_0)\cl$, $TC'_0$ petit. Si donc on a $f' \in TC' = C' \cap W'$, factorisons $f'$ en $p'i'$, avec
\[
i' \ \in \ (v_!(TC))\cl\ , \qquad p' \ \in \ (v^*)^{-1}(F)\ .
\tag*{(5.8.13)}
\]
\[
\begin{tikzcd}[column sep=small, row sep=small]
X' \arrow[rr, "f' \in C' \cap W'"] \arrow[dr, "i'"'] & & Y' \arrow[dl, dotted, shift right, "q'"'] \\
& X'_1 \arrow[ur, shift right, "p'"'] &
\end{tikzcd}
\tag*{$(*)$}
\]
Notons qu'on a
\msp{page 95}
\[
(v_!(TC))\cl \ \subset \ TC' \ (\ \defeq \ C' \cap W'\ )
\tag*{(5.8.14)}
\]
car $v_!(TC) \subset TC' \subset C'$, donc $(v_!(TC))\cl \subset \tilde{C}' = C'$, et il reste à prouver que $v_!(TC)\cl \subset W'$, i.e. $v^*((v_!(TC))\cl) \subset W$. Or je dis qu'on a même
\[
v^*((v_!(TC))\cl) \ \subset \ TC\ ,
\]
i.e.
\[
v_!(TC)\cl \ \subset \ (v^*)^{-1}(TC)\ .
\tag*{(5.8.15)}
\]
Cela résulte du fait que $TC$ est Q-clos, donc Q-saturé à gauche, ce qui implique, comme $v^*$ commute aux petites $\varinjlim$\,, que $(v^*)^{-1}(T)$ l'est également. D'autre part, il contient $v_!(TC)$, donc aussi le Q-saturé à gauche de $v_!(TC)$. Mais ce dernier est \textit{égal} à la Q-clôture $(v_!(TC))\cl$, par le théorème de factorisation, vu que $TC'_0 \subset v_!(TC) \subset (TC'_0)\cl$, avec $TC'_0$ petit. % [sic] (v^*)^{-1}(T) -> (v^*)^{-1}(TC), the set in question is TC

Ainsi, (5.8.13) et (5.8.14) impliquent $i' \in TC' = C' \cap W'$, donc $i' \in W'$, d'où $p' \in W'$, i.e. $v^*(p') \in W$, et comme par hypothèse
% foot: Version 9 février 2019
\printedpage{48}
% head: A. Grothendieck : Dérivateurs, Ch. XVII, Édition des Universités de Montpellier II et Paris VII
\msp{page 96}
(5.8.13) on a $v^*(p') \in F$, on a $v^*(p') \in TF$, i.e.
\[
p' \ \in \ (v^*)^{-1}(TF) \ = \ TF' \ (\ \defeq \ C'\ )\ , % [sic] C' -> C'_*, cf. (5.8.16) TF' ( = C'_* )
\]
l'égalité $(v^*)^{-1}(TF) = TF'$, i.e.
\[
(v^*)^{-1}(C_*) \ = \ (v_!(C))_*\ ,
\]
résultant de l'adjonction de $v_!$, $v^*$. Au total, on a donc
\[
p' \ \in \ TF' \ (\ = \ C'_*\ )\ ,
\tag*{(5.8.16)}
\]
et comme $f' \in C'$, le diagramme $(*)$, page 114, nous montre qu'il existe
\[
q' \ : \ Y' \longrightarrow X'_1\ ,
\]
avec $p'q' = \mathrm{id}$, $q'f' = i'$. Donc $f'$ est facteur direct de $i' \in v_!(TC)\cl$, donc on a $f' \in v_!(TC)\cl$, d'où la relation
\[
TC' \ \subset \ v_!(TC)\cl\ ,
\]
ce qui achève de prouver que $(W', C', F')$ est un triple de Quillen spécial et clos.

Les deux relations (5.8.4) ont été prouvées chemin faisant (la deuxième étant essentiellement triviale, et la première représentant la difficulté principale dans la démonstration). Reste à établir (5.8.5). Mais on sait déjà (5.8.9) que
\[
v_!(TC) \ \subset \ TC'\ ,
\]
\msp{page 97}
d'où
\[
v_!(TC)\cl \ \subset \ TC',
\]
puisque $TC'$ est Q-clos à gauche (ou parce que ça a été prouvé dans (5.8.14)), on a déjà prouvé l'inclusion en sens inverse (5.8.12 bis), assertion équivalente à (5.8.4) (et qui représentait la principale difficulté).

\textbf{Remarque 5.9.} Si on ne suppose pas que la structure sur $\mathcal{M}$ soit close, alors la démonstration donnée de 5.8 montre qu'on trouve pourtant sur $\mathcal{M}'$ une structure de Quillen spéciale, en définissant cette fois $C'$ comme dans 4.6. Mais je doute qu'on aura jamais à travailler avec des triples de Quillen spéciaux qui ne soient aussi clos. De toute façon, si $W$ est stable par petites $\varinjlim$ filtrantes\footnote[28]{ce qui implique la stabilité par facteurs directs, et implique aussi la même propriété pour $W' \defeq (v^*)^{-1}(W)$.} (et je doute qu'on trouve des triples de Quillen spéciaux en dehors de ce cas), alors si $(W, C, F)$ est un triple de Quillen spécial, il en est de même de $(W, \tilde{C} = C^\natural, F)$, et je présume qu'on a tout intérêt à remplacer le triple initial par celui-ci. % footnote mark printed in the text as superscript (28)
% foot: Version 9 février 2019

\end{document}
